%
% hub_display_reconciliation_no_preemption_buggy.tex
%
% FIDELITY CONTROL, not the design. Removes the socket-eviction half of
% single-owner preemption from HubIdentify: hubOwner still correctly
% replaces (I3 stays intact -- this variant isolates I1/I2 from I3), but a
% preempted predecessor's socket is never closed and its scenes are never
% orphaned. See docs/hub_display_reconciliation.tex, Section "Fidelity:
% reproducing the shipped defect", for the full trace and the invariant
% checks that must return FOUND against this variant (and NOT found
% against the correct spec).
%
% Type-check:  fuzz docs/hub_display_reconciliation_no_preemption_buggy.tex
% Model-check:
%   probcli docs/hub_display_reconciliation_no_preemption_buggy.tex \
%       -model_check -nodead -goal "..." -p DEFAULT_SETSIZE 3
%

\documentclass[a4paper,10pt,fleqn]{article}
\usepackage[margin=1in]{geometry}
\usepackage{fuzz}
\usepackage[colorlinks=true,linkcolor=blue,citecolor=blue,urlcolor=blue]{hyperref}

\begin{document}

\title{Hub/Display Scene Reconciliation --- Fidelity Control (No Preemption)}
\author{Buggy variant: HubIdentify never evicts a predecessor's socket}
\date{August 2026}
\maketitle

% ============================================================
\section{Introduction}
% ============================================================

This is the fidelity control for \texttt{docs/hub\_display\_reconciliation
.tex}, not the design.  It removes only the socket-eviction half of
single-owner preemption from \texttt{HubIdentify}: \texttt{hubOwner} still
correctly replaces (I3 stays intact, isolating this finding from a
multi-owner one), but a preempted predecessor's socket is never closed
and its scenes are never orphaned --- \texttt{SocketListener
.remove\_client} never runs.  Read the main specification first; this
document states only what differs, in Section~\ref{sec:identify}.

% ============================================================
\section{Basic Types}
% ============================================================

\begin{zed}
  [CONNECTION\_ID, SCENE\_ID]
\end{zed}

Unchanged from the main specification: three connections and two scenes,
bounded by ProB's \texttt{DEFAULT\_SETSIZE}.

% ============================================================
\section{Free Types}
% ============================================================

\begin{zed}
  OWNER ::= orphan | conn \ldata CONNECTION\_ID \rdata
\end{zed}

\begin{zed}
  FLAG ::= set | clear
\end{zed}

Unchanged from the main specification.

% ============================================================
\section{State}
% ============================================================

\begin{schema}{Recon}
  hubOwner : \power CONNECTION\_ID \\
  liveConns : \power CONNECTION\_ID \\
  sceneOwner : SCENE\_ID \pfun OWNER \\
  manifested : \power SCENE\_ID \\
  manifestDeclared : FLAG
\end{schema}

Unchanged from the main specification --- see there for the full
description of each field.

% ============================================================
\section{Initialization}
% ============================================================

\begin{schema}{Init}
  Recon'
\where
  hubOwner' = \emptyset \\
  liveConns' = \emptyset \\
  sceneOwner' = \emptyset \\
  manifested' = \emptyset \\
  manifestDeclared' = clear
\end{schema}

Unchanged from the main specification.

% ============================================================
\section{Operations}
% ============================================================

\subsection{HubIdentify --- BUGGY: no socket eviction}
\label{sec:identify}

\textbf{This is the one operation that differs from the design.}  A
connection \texttt{c?} declares \texttt{kind="hub"} and \texttt{hubOwner}
still correctly replaces --- I3 is deliberately left intact by this
variant, isolating the I1/I2 finding from a multi-owner finding.  But the
predecessor is never evicted from \texttt{liveConns}, and no scene is
orphaned: this models \texttt{SocketListener.remove\_client} never being
called, so the old socket stays open and \texttt{RenderLoop
.\_on\_client\_disconnected} never runs for it.  Compare against the
correct \texttt{HubIdentify} in \texttt{hub\_display\_reconciliation.tex}:
the correct version evicts \texttt{hubOwner}'s prior members from
\texttt{liveConns} and orphans their scenes in the same step; this one
does neither.

\begin{schema}{HubIdentify}
  \Delta Recon \\
  c? : CONNECTION\_ID
\where
  liveConns' = liveConns \cup \{ c? \} \\
  sceneOwner' = sceneOwner \\
  hubOwner' = \{ c? \} \\
  manifested' = manifested \\
  manifestDeclared' = clear
\end{schema}

\subsection{HubManifest, SceneInstall, ConnectionDrop --- unchanged}

Identical to the main specification: only \texttt{HubIdentify}'s
preemption step is removed.

\begin{schema}{HubManifest}
  \Delta Recon \\
  c? : CONNECTION\_ID \\
  m? : \power SCENE\_ID
\where
  c? \in hubOwner \\
  manifested' = m? \\
  sceneOwner' = \{ s : \dom sceneOwner | \\
  \quad~ s \in m? \lor sceneOwner~s = conn(c?) \} \dres sceneOwner \\
  hubOwner' = hubOwner \\
  liveConns' = liveConns \\
  manifestDeclared' = set
\end{schema}

\begin{schema}{SceneInstall}
  \Delta Recon \\
  c? : CONNECTION\_ID \\
  s? : SCENE\_ID
\where
  c? \in liveConns \\
  sceneOwner' = sceneOwner \oplus \{ s? \mapsto conn(c?) \} \\
  hubOwner' = hubOwner \\
  liveConns' = liveConns \\
  manifested' = manifested \\
  manifestDeclared' = manifestDeclared
\end{schema}

\begin{schema}{ConnectionDrop}
  \Delta Recon \\
  c? : CONNECTION\_ID
\where
  hubOwner' = hubOwner \setminus \{ c? \} \\
  liveConns' = liveConns \setminus \{ c? \} \\
  sceneOwner' = sceneOwner \oplus \\
  \quad~ (\{ s : \dom sceneOwner | sceneOwner~s = conn(c?) \} \cross \{ orphan \}) \\
  manifested' = manifested \\
  manifestDeclared' = clear
\end{schema}

% ============================================================
\section{Verification: the goals that must flip}
% ============================================================

Against this variant, I1's negation must be \emph{found} --- the trace in
\texttt{hub\_display\_reconciliation.tex}, Section ``Fidelity: reproducing
the shipped defect,'' reaches it directly:
\texttt{HubIdentify}(old); \texttt{SceneInstall}(old, s1);
\texttt{HubIdentify}(new); \texttt{HubManifest}(new, \{\});
\texttt{SceneInstall}(old, s1).
\begin{verbatim}
  # must be FOUND against this variant (NOT found against the correct spec)
  probcli docs/hub_display_reconciliation_no_preemption_buggy.tex \
    -model_check -nodead \
    -goal "manifestDeclared = set &
           (card({s | s : SCENE_ID & s : dom(sceneOwner) &
                 not(sceneOwner(s) = orphan) &
                 not(card(hubOwner /\ {c | c : CONNECTION_ID &
                       sceneOwner(s) = conn(c)}) > 0) }) > 0
           or
           card({s | s : SCENE_ID & s : dom(sceneOwner) &
                 sceneOwner(s) = orphan & not(s : manifested)}) > 0)" \
    -p DEFAULT_SETSIZE 3
\end{verbatim}

I3 must stay \emph{not found} --- this variant deliberately leaves
\texttt{hubOwner} correctly replacing, so it proves nothing about I3 on
its own and must not be mistaken for a multi-owner counterexample:
\begin{verbatim}
  # must remain NOT found -- this variant does not touch hubOwner's
  # replace-not-union behaviour, so I3 is not the finding here
  probcli docs/hub_display_reconciliation_no_preemption_buggy.tex \
    -model_check -nodead -goal "card(hubOwner) > 1" -p DEFAULT_SETSIZE 3
\end{verbatim}

\end{document}
